\input{preamble}

\begin{document}

\title{The Price of a Frozen Surrogate in Support-Aware Compression:\\
Separation, Starvation and the Cost of Repair}

\author{\IEEEauthorblockN{Vu Phan An Nguyen}
\IEEEauthorblockA{\textit{Rider University} \\
Lawrenceville, NJ, USA}}

\maketitle

\begin{abstract}
Which observations are worth buying depends on the alternatives to the decision, not only on
the observation model. We study a compression decision whose loss contains moving extrema,
and a surrogate that freezes the extremum cells while dropping their ordering constraints.
Both objectives optimise over the same designs and the same centred parameter class.
For a family with $r$ middle-cell pairs of variance $M$, the surrogate value is $9/512$,
whereas the true value lies between $9/[128(4+2rM)]$ and $1/(8rM)$.
Every surrogate-optimal design assigns zero mass to those pairs and has zero true information;
every $\varepsilon$-optimal design has true information at most $64\varepsilon/(9rM)$.
Thus the failure persists at unit variance as the number of neglected directions grows.
We quantify its repair. Sacrificing a fraction $\beta$ of surrogate value permits true
information between $9\beta/[128(4+2rM)]$ and $\beta/(8rM)$, and near $\beta=0$ the upper end
is attained to second order, so the exchange rate $1/(8rM)$ is exact.
A certificate based on the surrogate's dual margins extends this design tradeoff to a budget
law for adaptive policies, including unbounded stopping: at expected cost $N$, the average
design must pay $N\,\mathrm{sac}(\bar w)\ge8rM\log(1/\delta)(1-o(1))$.
On any channel the repair frontier $R(b)$ bounds the budgeted rate: a policy whose design
honours a sacrifice budget $b$ pays $\mathbb E|N|\,R(b)\ge(1-\delta)\log(1/\delta)-\log2$, also
for unbounded stopping. On Gaussian cells, over the same class of policies, a tracking learner
holding a design inside the budget attains the rate $1/R(b)$, which is therefore optimal; on two
Gaussian arms $R(b)=b(1-b)\Delta^2/2$ exactly for $b\le\frac12$.
An ideal learner attains the true characteristic time, while a specified learner that
re-solves the frozen surrogate at its estimate needs $\Omega(\log^2(1/\delta))$ samples under
square-root forcing. The distinction is an optimised failure of alternative geometry, followed
by a quantitative cost of correcting it. Lean~4 proofs support the stated model and policy
results. Exact rational certificates and paired simulations illustrate the mechanism;
they also show why neither a finite-precision learner nor a favourable finite panel can be
identified with the asymptotically optimal ideal algorithm.
\end{abstract}

\begin{IEEEkeywords}
compression, alternative sets, experimental design, surrogate objectives, sequential testing,
formal verification
\end{IEEEkeywords}

\section{Introduction}
\label{sec:intro}

\subsection{The question}

Let a world be $\theta$, let $T$ be a retained coordinate set, and let $L_{\theta,K}(T)$ be the
compression loss of $T$ on the deployment support $K$. A learner must return a $T$ whose
regret is at most $\varepsilon$, and it pays for observations. Its cost is governed by an
information game. The inner infimum of that game ranges over the \emph{alternatives}: the
worlds in which the returned $T$ would be wrong.

Computing that infimum requires the alternative set. The shortcut studied here freezes the cells
realising each maximum and minimum, and asks how far the observations must move to reverse
one linear contrast. We call this the \emph{frozen surrogate}. The observation model stays
correct, the stopping test stays valid and the decision is unchanged. Only the geometry used
to score designs is altered. This paper measures what that alteration costs, and what it
costs to undo.

\subsection{Results}

Five results carry the paper.
\begin{enumerate}
\item[1)] \textbf{Separation after optimisation} (Theorem~\ref{thm:sep}). Both games optimise
the allocation and range over the same centred class. Even so, the true game is a factor
$\Theta(rM)$ below the frozen one when $rM\ge1$, with absolute comparison constants; the factor is already present at unit variance. The
error has no fixed sign: on a second compression class the surrogate underprices by a factor
$M/96$, so no constant relates it to the truth.
\item[2)] \textbf{Starvation} (Theorem~\ref{thm:starve}). Every optimum of the surrogate puts
zero weight on the cells that carry the decision. An approximate optimum leaks in proportion
to its suboptimality.
\item[3)] \textbf{The price of repair} (Theorems~\ref{thm:frontier}, \ref{thm:sacrifice},
\ref{thm:sacstop}, \ref{thm:dilute}, \ref{thm:forcing}). A design that sacrifices a fraction
$\beta$ of surrogate value recovers $\Theta(\beta/rM)$ of true information when $rM\ge1$, and
the frontier's slope at $\beta=0$ is exactly $\frac{1}{8rM}$ for every $r$ and $M$. The displayed finite-parameter bounds remain valid for every $M>0$. Operationally, every
$\delta$-correct policy, including one that stops at an unbounded random time, must sacrifice a
budget of $8rM\log(1/\delta)$ query-units of surrogate value to first order. The rate-optimal
learner pays that budget up to a factor of at most $\frac{32}{9}+\frac{64}{9rM}$.
The budget can be spread over any run length: the relative sacrifice falls as $N_{0}/N$ while
the policy stays $\delta$-correct (Theorem~\ref{thm:dilute}).
A learner that meets its demand for informative queries only through forcing at rate
$t^{a-1}$ pays at least order $\log^{1/a}(1/\delta)$.
\item[4)] \textbf{Learners} (Theorems~\ref{thm:ideal}, \ref{thm:surlearner}). A completely
defined ideal learner attains the true rate. The learner that tracks the surrogate, at the
truth or re-solved at its own estimate, pays a whole extra power of $\log(1/\delta)$.
\item[5)] \textbf{The repair law} (Theorems~\ref{thm:repairlaw}, \ref{thm:optrate}). On any
channel, a policy whose realised design honours a sacrifice budget $b$ pays
$\mathbb E|N|\,R(b)\ge(1-\delta)\log(1/\delta)-\log2$, at a horizon and, for the expected
design of the whole run, at an unbounded stopping time. On Gaussian cells, over one class of
policies, the rate $1/R(b)$ is a lower bound and a tracking learner that holds a design inside
the budget attains it. So the repair frontier $R(b)$ is the optimal budgeted rate. The budget classes are nested, while an
expected-sacrifice constraint escapes them (Proposition~\ref{prop:classes}). On two Gaussian
arms the law is exact: $R(b)=b(1-b)\Delta^2/2$ (Proposition~\ref{prop:twoarm}).
\end{enumerate}
Three further facts are what make these statements about compression rather than about a
parameter class: the construction is an observation model, the alternative set is exactly a
union of polyhedra, and a margin certificate is not a decision (Section~\ref{sec:model}).

\subsection{Prior work, and what does not follow from it}

Information games with several correct answers, and alternative sets that decompose into
finitely many convex regions, are developed in the literature. Degenne and
Koolen~\cite{degenne2019} treat multiple correct answers, with a lower bound and an algorithm
that attains it. The characteristic time and Track-and-Stop are due to Garivier and
Kaufmann~\cite{garivier2016}, the change-of-measure inequality to Kaufmann, Capp\'e and
Garivier~\cite{kaufmann2016}, and Jourdan et al.~\cite{jourdan2021} separate the answer set from
the action set in combinatorial pure exploration. We inherit that machinery and do not claim
it. Convex relaxations of the alternative geometry, used to make such algorithms tractable,
appear in~\cite{shukla2024}; the surrogate studied here is not a relaxation, and it misprices in
both directions. Wang and You~\cite{wangyou2026} study ranking-and-selection with multiple
correct answers and non-answerable estimates. That an information proxy can select a design
lacking identifiability, at every one of its optima under a stated condition, is
known~\cite{overstall2022}. So is the fact that acquisition should target the decision rather
than the predictor~\cite{filstroff2021}, error amplification under
misspecification~\cite{tang2026}, and decision-focused sequential design with a failing
prediction-focused instance~\cite{wan2026}. When a proxy suffices for a decision is studied
directly~\cite{schutte2026}. Bennouna, Bennouna, Amin and Ozdaglar characterise exactly when a
noiseless linear dataset, together with prior knowledge and the task's uncertainty set,
determines an optimal decision~\cite{bennouna2025,bennouna2026}. Their geometric condition
identifies which cost directions must be captured; it is not an unconditional full-rank
requirement. We ask a different question: what those directions cost to
measure under noise, and what a surrogate that drops them costs to repair.

Pricing identification in a second currency is also known. Yang, Tan and Jin~\cite{yang2024}
bound the regret incurred before a best-arm identification stops. Huang~\cite{huang2026}
certifies representation adequacy at minimum incurred loss, with a cost-weighted
characteristic time and a matching Track-and-Stop for a fixed observation kernel, and leaves
representation repair open. Repair there means changing the representation, its kernels or its
loss after certification; we do not address it. Here the observation model is fixed and only the
allocation is repaired. The budget is paid in the surrogate's own information game, which is
concave in the design rather than a per-query cost, and the results that price it are the
geometry of this surrogate: the frontier, its exact slope and the sacrifice law.

The generic statements are therefore taken. What remains is narrower. The model is correct,
the test is valid and the decision is unchanged, and still the frozen geometry changes what a
design can learn and what a learner pays. Table~\ref{tab:prior} states the difference at the
level of theorems.

\begin{table}[h]
\caption{Closest research threads and the specific quantitative results compared here.}
\label{tab:prior}
\centering
\footnotesize
\begin{tabular}{@{}p{0.30\columnwidth}p{0.62\columnwidth}@{}}
\toprule
Nearest source & Specific result compared here \\
\midrule
Proxy insufficiency~\cite{schutte2026,overstall2022} & The bracket
$\frac{4}{4+2rM}\le\Gamma_{\rm true}/\Gamma_{\rm fr}\le\frac{64}{9rM}$ after optimisation, and
the leak $J_{\rm true}\le\frac{64\varepsilon}{9rM}$ of every $\varepsilon$-optimal design
(Thms.~\ref{thm:sep}, \ref{thm:starve}); the every-optimum form alone is known \\
Data sufficiency~\cite{bennouna2025} & The noisy price of the directions a surrogate drops: the
frontier and the budget (Thms.~\ref{thm:frontier}, \ref{thm:sacrifice}) \\
Cost-weighted certification~\cite{huang2026,yang2024} & A budget in the surrogate's concave
value, and its repair price; spreading it (Thms.~\ref{thm:sacrifice}, \ref{thm:dilute}) \\
Multiple-answer track-and-stop~\cite{degenne2019} & The repair frontier
$\Theta(\beta/rM)$ for $rM\ge1$ and the sacrifice budget $N\beta\ge8rM\log(1/\delta)(1-o(1))$
(Thms.~\ref{thm:frontier}, \ref{thm:sacrifice}) \\
Characteristic time, multiple answers~\cite{garivier2016,degenne2019} & The budgeted optimal rate
$1/R(b)$ over the nested budget classes, and its attainment by a fixed-design tracking learner
(Thm.~\ref{thm:repairlaw}, Thm.~\ref{thm:optrate}) \\
Decision-focused design~\cite{wan2026} & The $\Omega(\log^{2}(1/\delta))$ cost of the
learner that re-solves the surrogate at its estimate, against the ideal learner's
$\Theta(\log(1/\delta))$ (Thms.~\ref{thm:ideal}, \ref{thm:surlearner}) \\
\bottomrule
\end{tabular}
\end{table}

\subsection{What is proved, and how}

Three evidence classes appear, and they are never interchanged.
\begin{itemize}
\item \PROVED: a theorem stated at the paper's objects and checked in
Lean~4~\cite{demoura2021} against Mathlib~\cite{mathlib2020}.
\item \CERT: an exact rational certificate.
\item \NUM: a floating-point or Monte Carlo result.
\end{itemize}
Each formal theorem carries a line naming its Lean declaration. Elementary consequences
assembled here carry an explicit proof and are distinguished from separately formalized wrappers. The supplement gives the proof
architecture, every supporting declaration, and the status of results superseded during
preparation. The checked proof modules use no \texttt{sorry} or \texttt{native\_decide};
their audited axioms lie within \{\texttt{propext}, \texttt{Classical.choice}, \texttt{Quot.sound}\}.
Every hypothesis of a headline is either a model assumption displayed in it or discharged.

\section{The model, the decision and its alternatives}
\label{sec:model}

\paragraph{The loss} Parametrise the world linearly, $F_{\theta}(x)=\phi(x)^{\top}\theta$ on a
finite domain with a fixed reference and a deployment support $K$. For a retained set $T$ the
residual at a support point is a linear form $a_{T}(x)^{\top}\theta$, and
\begin{equation}
L_{\theta,K}(T)=\tfrac12\Bigl[\max_{x\in K}a_{T}(x)^{\top}\theta
-\min_{x\in K}a_{T}(x)^{\top}\theta\Bigr].
\label{eq:loss}
\end{equation}
Query $e$ returns $Y_{e}\sim\mathcal N(\theta_{e},\nu_{e})$, and a design is an allocation
$w\in\Delta$. For a set of alternatives $\mathcal A$, the value of a design and the game are
\begin{equation}
J(w)=\inf_{\lambda\in\mathcal A}\ \tfrac12\sum_{e}\frac{w_{e}(\lambda_{e}-\theta_{e})^{2}}{\nu_{e}},
\qquad \Gamma=\sup_{w\in\Delta}J(w).
\label{eq:game}
\end{equation}
The summand is the per-query Gaussian divergence. Thus $1/\Gamma$ is the candidate
characteristic time; its lower-bound interpretation and attainability require the policy
results proved below. This definition alone does not establish an algorithmic rate.

\leanbox{\lean{residual\_eq\_inner}, \lean{loss\_eq\_half\_range},
\lean{selectedRadius\_eq\_compressionLoss} in \lean{LeanPaperAModelBridgeV1};
\lean{energy\_eq\_sum\_gaussKL} in \lean{LeanPaperAStrictV1}}

\paragraph{The construction} Two arms, $A$ with cells $a_{0},a_{1}$ and $B$ with anchors
$b_{2},b_{3}$ and $2r$ middle cells grouped in $r$ pairs. The response is additive,
$F_{\lambda}(a,b)=\lambda_{A}(a)+\lambda_{B}(b)$. Anchors and arm-$A$ cells have variance $1$
and middle cells variance $M>0$. The world $\theta_{C}$ puts arm $A$ at $(-\frac12,\frac12)$,
the anchors at $(-\frac18,\frac18)$ and every middle cell at $0$. The class $\Theta_{0}$
requires each arm's cells to sum to zero.

\begin{theorem}[The construction is an observation model]
\label{thm:bridge}
$\lambda\in\Theta_{0}$ if and only if both uniform-product marginals recover their coordinates.
On $\Theta_{0}$, $L_{\lambda}(\{A\})=\frac12\mathrm{range}(\lambda_{B})$ and
$L_{\lambda}(\{B\})=\frac12\mathrm{range}(\lambda_{A})$. Retaining $A$ is $\varepsilon$-unsafe
if and only if $\mathrm{range}(\lambda_{B})>\mathrm{range}(\lambda_{A})+2\varepsilon$, and this
set is exactly the union over pairs $(p,q)$ of arm-$B$ cells of the polyhedra
$\{\lambda_{1}-\lambda_{0}+2\varepsilon<\lambda_{p}-\lambda_{q},\
\lambda_{0}-\lambda_{1}+2\varepsilon<\lambda_{p}-\lambda_{q}\}$.
\end{theorem}

\leanbox{\lean{centred\_iff\_marginals}, \lean{lossRetainA\_eq}, \lean{unsafeA\_iff\_range\_gap}
in \lean{LeanPaperACentredModelV1}; \lean{unsafe\_iff\_union\_of\_polyhedra} in
\lean{LeanPaperARegionV1}}

The unsafe condition is strict. Its closure contains ties, which are not unsafe. The game the
decision defines is therefore taken over the \emph{strict} alternatives
$\mathcal A_{\rm true}=\{\mathrm{range}(\lambda_{A})<\mathrm{range}(\lambda_{B})\}\cap\Theta_{0}$,
and every bound below is approached from inside that strict set.

\paragraph{The two games} The true game uses $\mathcal A_{\rm true}$. The frozen surrogate
keeps arm $A$'s cells and arm $B$'s anchors as the extrema and tests the one contrast
$(\lambda_{a_{1}}-\lambda_{a_{0}})-(\lambda_{b_{3}}-\lambda_{b_{2}})\le0$ over $\Theta_{0}$. The
two sets are not nested. The surrogate admits parameters that break the ordering, and it
misses alternatives that move a middle cell past an anchor.

\paragraph{A margin certificate is not a decision} Suppose the certificate condition implies
that $T$ is safe. Then the certificate's alternatives contain the decision's, so certifying is
never cheaper than deciding. On $F=ax_{1}+bx_{2}+cx_{1}x_{2}$ the decision is the sign of $a-b$,
while the usual margin certificate asks $a-b\ge2c$. At $a=3$, $b=1$, $\varepsilon=0.1$ the
certificate costs $19.4$ times the decision at $c=0.25$ and $1089$ times at $c=0.9$. It is
unattainable for $c\ge1$, where the decision is still settled. \NUM{}

\leanbox{\lean{certification\_costs\_at\_least\_as\_much} in
\lean{LeanPaperAOracleCertificateV1}}

\section{Separation after the allocation is optimised}
\label{sec:separation}

\begin{theorem}[Separation]
\label{thm:sep}
For every $r\ge1$ and $M>0$,
\[
\Gamma_{\rm fr}=\frac{9}{512},\qquad
\frac{9}{128(4+2rM)}\ \le\ \Gamma_{\rm true}\ \le\ \frac{1}{8Mr},
\]
and hence
\[
\frac{4}{4+2rM}\ \le\ \frac{\Gamma_{\rm true}}{\Gamma_{\rm fr}}\ \le\ \frac{64}{9rM}.
\]
The same bounds hold for the game over the strict alternatives. There, the Gaussian divergence
of the observation law is the energy of~\eqref{eq:game} for every centred candidate.
\end{theorem}

\leanbox{\lean{frozenGameC₀\_eq}, \lean{trueGameC₀\_le}, \lean{trueGameC₀\_ge},
\lean{centred\_separation} in \lean{LeanPaperACentredV1}; \lean{strict\_game\_bounds},
\lean{actualKLValue\_eq} in \lean{LeanPaperAEndgameV1}; \lean{gaussianKLIdentification\_holds}
in \lean{LeanPaperAGaussianKLV1}}

\emph{Why.} The upper bound is a pigeonhole. Some pair of middle cells carries at most $1/r$
of any design, and moving that pair apart by $\frac12+d$ reverses the decision. The move is a
strict alternative inside $\Theta_{0}$ for every $d>0$. The lower bound is the
variance-proportional design $w_{e}=\nu_{e}/(4+2rM)$, at which the energy is the plain squared
distance. Reversing the decision forces a contrast of at least $g=\frac34$ across four cells,
and Cauchy--Schwarz turns that into $\frac{9}{64}$ of squared distance. The same $g$ gives
$\Gamma_{\rm fr}=g^{2}/32$. Both infima are restricted to $\Theta_{0}$, and restricting an
infimum can only raise it. So neither bound is inherited from a larger set: each witness is
proved to lie in the class.

\emph{The surrogate can also underprice.} The centred construction shows one direction. A second
compression class shows both, again after the sup over designs. It has two relations of four
cells, the decision compares their ranges, every table is the response table of an additive
world, and one cheap cell has variance $M$. On one centred world of this class
$\Gamma_{\rm true}/\Gamma_{\rm fr}\le64/M$, and on another $\Gamma_{\rm true}/\Gamma_{\rm fr}\ge M/96$.
Both games exist, and no constant factor relates the surrogate to the truth on the class.

\leanbox{\lean{low\_separation\_at\_world}, \lean{high\_separation\_at\_world},
\lean{exists\_low\_games}, \lean{exists\_high\_games}, \lean{no\_constant\_factor\_on\_families},
\lean{additive\_response\_shift} in \lean{LeanPaperASeparationV1}}

\section{The surrogate misdesigns}
\label{sec:starve}

A ratio of optimal values can diverge while both objectives select the same design. What
matters is the design the surrogate selects.

\begin{theorem}[Starvation]
\label{thm:starve}
For every allocation $w$, $J_{\rm fr}(w)\le\frac{9}{512}\,a(w)$, where $a(w)$ is the mass on
arm $A$ and the anchors. Consequently:
\begin{enumerate}
\item[(i)] every $w$ with $J_{\rm fr}(w)=\Gamma_{\rm fr}$ puts zero mass on the middle cells,
and $J_{\rm true}(w)=0$;
\item[(ii)] every $w$ with $J_{\rm fr}(w)\ge\Gamma_{\rm fr}-\varepsilon$ puts at most
$\frac{512}{9}\varepsilon$ on the middle cells, and $J_{\rm true}(w)\le\frac{64\varepsilon}{9rM}$.
\end{enumerate}
\end{theorem}

\leanbox{\lean{frozenValueC₀\_le}, \lean{starvation\_at\_optimum₀}, \lean{near\_optimal\_starvation₀},
\lean{near\_optimal\_trueValueC₀\_le} in \lean{LeanPaperACentredV1};
\lean{frozen\_optimum\_trueValueStrict\_zero} in \lean{LeanPaperAStrictV1}}

Solving the surrogate more accurately therefore drives the design towards one that cannot
distinguish the alternatives the compression turns on. Part~(ii) is what a learner running on
the surrogate actually has: an approximate optimum, which leaks in proportion to its own
suboptimality.

\section{The price of repair}
\label{sec:repair}

Write the \emph{relative sacrifice} of a design as
$\mathrm{sac}(w)=1-J_{\rm fr}(w)/\Gamma_{\rm fr}$. It is a diagnostic of a design, evaluated at
the true world, and it never enters a learner.

\subsection{At a design}

\begin{theorem}[The repair frontier]
\label{thm:frontier}
For every $r\ge1$, $M>0$ and $\beta\in[0,1]$, the best true information of a design with
$\mathrm{sac}(w)\le\beta$ satisfies
\[
\frac{9\beta}{128(4+2rM)}\ \le\ \sup_{\mathrm{sac}(w)\le\beta}J_{\rm true}(w)\ \le\
\frac{\beta}{8rM}.
\]
This holds on the strict game as well. Pointwise, every allocation has
$J_{\rm true}(w)\le\mathrm{sac}(w)/(8rM)$. The design
$w_{\beta}=(1-\beta)w_{\rm unif}+\beta w_{\rm cert}$ achieves at least the displayed lower bound.
Moreover, for $0<\beta\le\frac12$,
\[
\frac{\beta}{8rM+128\beta}\ \le\ \sup_{\mathrm{sac}(w)\le\beta}J_{\rm true}(w)\ \le\
\frac{\beta}{8rM},
\]
and the two ends differ by at most $2\beta^{2}/(rM)^{2}$. So the frontier's slope at $\beta=0$ is
exactly $\frac{1}{8rM}$. We do not assert an exact formula for the frontier at positive $\beta$.
\end{theorem}

\leanbox{\lean{value\_concave}, \lean{frontier\_two\_sided} in \lean{LeanPaperAFrontierV1};
\lean{exchange\_strict}, \lean{strict\_frontier} in \lean{LeanPaperASacrificeV1};
\lean{vB\_isAlloc}, \lean{vB\_frozen}, \lean{trueValue\_vB\_ge}, \lean{frontier\_slope},
\lean{slope\_gap} in \lean{LeanPaperARepairSlopeV1}; \lean{trueValueC₀\_le\_trueValueStrict} in
\lean{LeanPaperAStrictV1}}

Each game's value is an infimum of functions linear in the design, hence concave in it, so
mixing two designs loses at most linearly. The upper side is Theorem~\ref{thm:starve}(ii) with
$\varepsilon=\beta\Gamma_{\rm fr}$. The exchange rate is linear in the sacrifice, and it
carries the same $\frac{1}{rM}$ as the separation. The sharp lower side mixes the uniform-anchor
design, with weight $1-\beta$, and the uniform design on the $2r$ middle cells, with weight
$\beta$. Against a strict
alternative, the four frozen cells and the moved middle pair each must absorb part of a contrast;
a weighted Cauchy--Schwarz over the cells that move, split by which cells realise arm $B$'s
extrema, gives $\beta/(8rM+128\beta)$ in every case. The certificate constant $8rM$ of the
sacrifice law below is therefore the exact first-order exchange rate, not an artefact of the
bound.

\subsection{A certificate that extends beyond this family}
\label{sec:repaircertificate}
The pair construction makes the constants explicit. The inference from a flawed design
criterion to a repair budget uses less structure. Let $J_f,J_t$ be nonnegative concave design
values on a finite simplex, let $\Gamma_f=\max_wJ_f(w)>0$, and write
$s(w)=1-J_f(w)/\Gamma_f$. Suppose a frozen certificate $c$ and true-alternative load vectors
$\ell_j$ satisfy, for every design $w$,
\[
J_f(w)\le\sum_e w_ec_e,\qquad
J_t(w)\le\sum_e w_e\ell_j(e)\quad\text{for every }j.
\]
A distribution over frozen alternatives supplies $c$; each true alternative supplies one
load. These are verifiable information inequalities, not estimates of the two optima.

\begin{proposition}[Certificate for repair cost]
\label{prop:repaircertificate}
Suppose $\nu_j\ge0$, $\Lambda=\sum_j\nu_j>0$, and
\begin{equation}
\sum_j\nu_j\ell_j(e)\le1-\frac{c_e}{\Gamma_f}\quad\text{for every cell }e.
\label{eq:repaircover}
\end{equation}
Then $\Lambda J_t(w)\le s(w)$ for every design. Consequently
$R(b):=\sup_{s(w)\le b}J_t(w)\le b/\Lambda$.
If expected counts $N_e\ge0$ have $N=\sum_eN_e>0$ and meet
$\sum_eN_e\ell_j(e)\ge D$ for every $j$, then
\begin{equation}
N s(N_{\cdot}/N)\ge\Lambda D.
\label{eq:generalrepairbudget}
\end{equation}
Conversely, if $w_*$ attains $\Gamma_f$, and $v$ has $s(v)>0$, then for every $b\ge0$,
\begin{equation}
R(b)\ge\min\{1,b/s(v)\}\,J_t(v).
\label{eq:clippedrepair}
\end{equation}
\end{proposition}
\begin{proof}
Multiply each true-load inequality by $\nu_j$, sum, and use~\eqref{eq:repaircover}:
\[
\Lambda J_t(w)\le\sum_e w_e\sum_j\nu_j\ell_j(e)
\le1-\frac{\sum_e w_ec_e}{\Gamma_f}\le s(w).
\]
The same calculation with $N_e$ in place of $w_e$ proves~\eqref{eq:generalrepairbudget}.
For the converse take $\lambda=\min\{1,b/s(v)\}$ and
$w=(1-\lambda)w_*+\lambda v$. Concavity gives $s(w)\le\lambda s(v)\le b$ and
$J_t(w)\ge\lambda J_t(v)$. Clipping $\lambda$ at one is necessary: extrapolation is not a design.
\end{proof}
\leanbox{\lean{exchange}, \lean{frontier\_upper}, \lean{budget}, \lean{weak\_duality},
\lean{repair\_mix}, \lean{repair\_clipped} in \lean{LeanPaperASacrificeDualityV1}.
Equation~\eqref{eq:clippedrepair} is \lean{repair\_clipped}.}

For our family, $c_e=\Gamma_f$ on the four frozen cells and zero on the middles. A strict
alternative moves one middle pair by $1/2+d$, and has load $(1/2+d)^2/(2M)$ on each of its
two cells. Set $\nu_j=8M/(1+2d)^2$. The cover is exact on every middle cell and zero on the
others; hence $\Lambda=8rM/(1+2d)^2$. Taking $d\downarrow0$ recovers both the frontier upper
bound and the policy budget $8rMD$ without replacing strict alternatives by unattained ties.
This is a sufficient certificate and a weak-duality bound, not a claim that every game's
repair frontier has this form or that these bounds are always tight.
\leanbox{\lean{cStar\_cert}, \lean{pairLoad\_cert}, \lean{nuStar\_cover},
\lean{construction\_frontier}, \lean{construction\_budget},
\lean{certificate\_pair\_satisfiable} in \lean{LeanPaperASacrificeDualityV1}}

\subsection{For a policy}

A design frontier does not say what a learner must give up. That requires every policy, its
own random sample sizes, and the correctness it must reach. Fix a randomised adaptive policy
with a finite horizon, whose answer event has probability at least $1-\delta$ at $\theta_{C}$
and at most $\delta$ at every strict alternative. Let $N_{e}$ be its expected number of queries
of cell $e$ at the truth, $N=\sum_{e}N_{e}$, and $\bar w=N_{\cdot}/N$ its average design. Put
$D=(1-\delta)\log(1/\delta)-\log2$, and take $\delta$ small enough that $D>0$; every
$\delta<\frac18$ qualifies (\lean{hD\_of\_small}). For larger $\delta$ the bounds below are void.

\begin{theorem}[The sacrifice law]
\label{thm:sacrifice}
Every such policy queries each middle pair at least $8MD$ times in expectation, and the middle
cells at least $8rMD$ times in total. Each middle query is paid for in surrogate value: the
middle mass of any allocation is at most its relative sacrifice. Hence
\[
N\cdot\mathrm{sac}(\bar w)\ \ge\ 8\,r\,M\,D .
\]
The hypotheses are satisfiable. The ideal learner of Theorem~\ref{thm:ideal}, truncated at a
Markov horizon, satisfies them at level $2\delta$. For that learner
$N\,\mathrm{sac}(\bar w)\le N\le\mathbb{E}\tau_{\delta}$, and
$1/\Gamma_{\rm true}\le\frac{128(4+2rM)}{9}$.
\end{theorem}

\leanbox{\lean{pair\_demand}, \lean{middle\_demand}, \lean{middle\_le\_sacrifice},
\lean{sacrifice\_cost\_law}, \lean{upper\_constant}, \lean{constant\_gap} in
\lean{LeanPaperASacrificeV1}; \lean{ideal\_sacrifice}, \lean{sacrifice\_hypotheses\_satisfiable}
in \lean{LeanPaperASacrificeIdealV1}}

\emph{What this says.} To make the true decision, a policy must give up a \emph{budget} of
surrogate value, $8rM\log(1/\delta)$ query-units to first order. It is not a fraction. A policy
may spend the budget slowly, with a small sacrifice over a long run, or quickly, but it cannot
avoid it. The rate-optimal learner spends at most $\frac{128(4+2rM)}{9}\log(1/\delta)$ to first
order. For $rM\ge1$, the budget is $\Theta(rM\log(1/\delta))$ with absolute comparison constants.
For arbitrary $r,M$ the exact bounds above are the statement; their ratio is
$\frac{32}{9}+\frac{64}{9rM}$.

\begin{theorem}[Spreading the budget]
\label{thm:dilute}
Let $\pi_{0}$ be a policy as above with $N_{0}$ expected real queries and average design
$\bar w_{0}$. Run it for its $n$ rounds, then query the four frozen cells (the two $A$ cells and the two $B$ anchors) round-robin for $4k$ more
rounds and answer from the first $n$. For every $k$, this policy is $\delta$-correct with the same
answer event. It has $N=N_{0}+4k$ expected real queries, it obeys the sacrifice law, and it spends
no more budget than $\pi_{0}$:
\begin{gather*}
8rMD\ \le\ N\,\mathrm{sac}(\bar w)\ \le\ N_{0}\,\mathrm{sac}(\bar w_{0}),\\
\mathrm{sac}(\bar w)\ \le\ \frac{N_{0}}{N_{0}+4k}.
\end{gather*}
\end{theorem}

\leanbox{\lean{sacrifice\_mix\_le}, \lean{padPol}, \lean{pad\_real}, \lean{lintegral\_cnt\_pad},
\lean{padCnt\_A}, \lean{dilution}, \lean{dilution\_small} in \lean{LeanPaperASacrificeDiluteV1}}

With $\pi_{0}$ the truncated ideal learner, $N_{0}\le\mathbb{E}\tau_{\delta}
=O((4+2rM)\log(1/\delta))$. So for every run length $N\ge N_{0}$, in steps of four, a
$\delta$-correct policy gives up a relative surrogate value of at most $N_{0}/N$. The budget
$N\,\mathrm{sac}(\bar w)$ stays between $8rMD$ and $N_{0}$. The padding is surrogate-optimal, and
concavity of the frozen value turns the mixture of designs into
$\mathrm{sac}\le\beta\,\mathrm{sac}(\bar w_{0})$ with $\beta=N_{0}/N$.

\begin{theorem}[The law for unbounded stopping]
\label{thm:sacstop}
Let a randomised adaptive policy have a measurable stopping rule that stops almost surely at
$\theta_{C}$, with no bound on the stopping time. Suppose it answers $A$ with probability at least
$1-\delta$ at $\theta_{C}$ and at most $\delta$ at every strict alternative, on its stopped runs.
Let $N_{e}$ be its expected number of queries of cell $e$ over the whole run, assumed finite. Then
$N\,\mathrm{sac}(\bar w)\ge8rMD'$, with $D'=(1-2\delta)\log\frac1{2\delta}-\log2>0$ (every
$\delta<\frac1{16}$ qualifies). A finite-horizon policy that satisfies Theorem~\ref{thm:sacrifice}
is such a policy.
\end{theorem}

\leanbox{\lean{outT\_unfired}, \lean{trunc\_A}, \lean{expCntSt}, \lean{sacrifice\_law\_stopped},
\lean{correctA\_ofFinite} in \lean{LeanPaperASacrificeStoppedV1}}

An unstopped run answers ``not $A$'', so no uniform stopping over the alternatives is needed.
Truncating at a horizon where at most $\delta$ of the truth's runs are still running gives a
finite-horizon policy that is correct at level $2\delta$. Its counts are at most the unbounded
run's, and the middle mass of any design is at most its sacrifice.

\emph{Why.} Each middle pair is moved by a strict alternative that only that pair's queries can
detect. The change of measure at the shift $\frac12+d$, with $d\to0$, gives the per-pair demand.
The surrogate's value is at most $\Gamma_{\rm fr}$ times the non-middle mass, which gives the
exchange.

\subsection{Forcing schedules}

A learner that follows the surrogate meets the demand for middle queries only through its
exploration. The sacrifice law then fixes its cost, whatever the schedule.

\begin{theorem}[The forcing law]
\label{thm:forcing}
For a policy as above, now with $D\ge0$, let $T$ be the number of real queries and $\mathrm{Mid}$
the number of middle queries at the horizon. Suppose $\mathrm{Mid}\le\varphi(T)+B$ on every
history, where
$\varphi$ is continuous and concave on $[0,\infty)$ and $\mathbb{E}B\le b$. Then
$\varphi(\mathbb{E}T)\ge8rMD-b$. For $\varphi(t)=ct^{a}$ with $0<a\le1$ and
$b\le 8rMD$,
\[
\mathbb{E}T\ \ge\ \Bigl(\frac{8rMD-b}{c}\Bigr)^{1/a}.
\]
\end{theorem}

\leanbox{\lean{forcing\_law}, \lean{power\_forcing\_law} in \lean{LeanPaperASacrificeV1};
\lean{forcing\_law\_stopped}, \lean{power\_forcing\_law\_stopped} in
\lean{LeanPaperAForcingStopV1}}

The law holds for unbounded stopping as well. If a policy stops almost surely at $\theta_C$ and is
correct on its stopped runs, if every truncation has $\mathrm{Mid}\le\varphi(T)+B_n$ with
$\mathbb EB_n\le b$, and if $\varphi$ is also nondecreasing, then
$\varphi(\bar T)\ge8rMD'-b$ for every bound $\bar T$ on the expected real queries of the whole
run, with $D'=(1-2\delta)\log\frac1{2\delta}-\log2\ge0$. The truncation at a horizon where at most
$\delta$ of the runs are unstopped is correct at level $2\delta$, and its expected real queries are
at most $\bar T$.

At $a=\frac12$, this gives a $\Omega(\log^{2}(1/\delta))$ lower bound. At $a=1$, the
same argument gives only a $\Omega(\log(1/\delta))$ lower bound; it does not by itself
prove an upper bound. Persistent corrective exploration can place positive mass outside the
surrogate's support. Whether a particular learner satisfies the pathwise count bound is a
property of that learner. Section~\ref{sec:learners} proves it for the two surrogate learners.

\section{Learners}
\label{sec:learners}

\subsection{The rule that the theorem actually analyses}
Let $\mathcal A_T$ be the strict centred alternatives to answer $T$, and let $N_e(t)$ and
$\widehat\theta_e(t)$ be the count and raw sample mean after $t$ observations. The variance
$\nu_e>0$ is known. For $a>0$ the stopping statistic and threshold are
\begin{align}
Z_t^a(T)&=\inf_{\lambda\in\mathcal A_T}
 \frac12\sum_e\frac{N_e(t)^2}{(a+N_e(t))\nu_e}
 (\widehat\theta_e(t)-\lambda_e)^2,\label{eq:ridgestop}\\
b_t(\delta)&=\log(1/\delta)+\frac12\sum_e\log(1+N_e(t)/a).
\label{eq:idealboundary}
\end{align}
This is the exact infimum of a ridge-weighted statistic. It is not the ordinary Gaussian
log-GLR, whose coefficient would be $N_e(t)/\nu_e$. The finite-count shrinkage is part of
the confidence argument, even though $N_e/(a+N_e)\to1$ under forcing.
The stopping time is
$\tau_\delta=\inf\{t\ge K:\max_T Z_t^a(T)>b_t(\delta)\}$, with value $+\infty$ when this set
is empty; the output is the first maximiser of $Z_t^a$ at the stop.
\leanbox{\lean{qAbs}, \lean{qAbs\_eq}, \lean{glr}, \lean{measurable\_glr} in
\lean{LeanPaperA7GLRMeasV1}; \lean{Zh}, \lean{stopAt}, \lean{answerAt}, \lean{tauI} in
\lean{LeanPaperA7IdealStopV1}; \lean{bthr} in \lean{LeanPaperA7ConfidenceV1}}

The query rule also reads only the history. First query each of the $K=4+2r$ cells once.
Thereafter let
\begin{gather*}
\Delta_{s+1}=\{w\in\Delta:(s+1)w_e\in\mathbb N\text{ for all }e\},\\
(T_s,w_s)\in\arg\max_{T,\,w\in\Delta_{s+1}}J_T(w;\widehat\theta(s)),
\end{gather*}
with a fixed tie order, where $J_T$ uses the entire alternative set, without a box.
Put $u_e=1/K$, $\rho_s=1/\sqrt{s+1}$, and $v_s=(1-\rho_s)w_s+\rho_su$.
After $t\ge K$ observations choose
\[
q_{t+1}\in\arg\max_e\left\{\sum_{s=K}^{t}v_{s,e}-N_e(t)\right\},
\]
again with fixed ties. The quota starts at zero after the warm start; the counts still
include the warm samples. This convention is the one in the formal tracking lemma.
\leanbox{\lean{netSet} in \lean{LeanPaperA7NetOracleV1}; \lean{oracle}, \lean{target},
\lean{quota}, \lean{pick}, \lean{measurable\_pick} in \lean{LeanPaperA7IdealLearnerV1};
\lean{famSpec} in \lean{LeanPaperA7IdealFamilyV1}}

Why does the threshold certify an answer? At a candidate true world $\theta$, define
\[
M_t(\theta)=\prod_e\sqrt{\frac{a}{a+N_e(t)}}
\exp\left\{\frac{N_e(t)^2(\widehat\theta_e(t)-\theta_e)^2}
{2(a+N_e(t))\nu_e}\right\}.
\]
It is a nonnegative martingale under the adaptive Gaussian observation law, starting at one.
If an incorrect answer $T$ crosses~\eqref{eq:idealboundary}, then
$\theta\in\mathcal A_T$ and $M_t(\theta)>1/\delta$.
Ville's inequality supplies correctness. Tracking, vanishing grid error and the direct
information-accumulation bound supply termination and the expected rate; correctness alone
would supply neither.

\begin{theorem}[The ideal learner]
\label{thm:ideal}
There is a completely defined learner with the following steps. It starts warm, with one query
per cell. At round $s$ it selects an answer and a design by an exact maximiser of $J$ over a
grid of mesh $\frac{1}{s+1}$. It mixes in the uniform design with weight
$\frac{1}{\sqrt{s+1}}$ and tracks the result. It stops at the ridge-statistic crossing
\eqref{eq:ridgestop}--\eqref{eq:idealboundary}. On the construction,
for every $a>0$:
\begin{itemize}
\item it stops with a wrong answer with probability at most $\delta$;
\item for every $\eta\in(0,\Gamma_{\rm true})$ there is $A_{\eta}$ such that, for every
$\delta\in(0,1)$, $\mathbb{E}\tau_{\delta}\le\frac{\log(1/\delta)}{\Gamma_{\rm true}-\eta}+A_{\eta}$;
\item every finite-horizon policy that is $\delta$-correct against the strict alternatives
makes at least $\frac{(1-\delta)\log(1/\delta)-\log2}{\Gamma_{\rm true}}$ real queries in
expectation, and the stopped learner meets this bound at level $2\delta$.
\end{itemize}
Hence $\mathbb{E}\tau_{\delta}/\log(1/\delta)\to1/\Gamma_{\rm true}$, where $\Gamma_{\rm true}$ is
the Gaussian-divergence game over the strict alternatives.
\end{theorem}

\leanbox{\lean{ideal\_family\_finish}, \lean{ideal\_family\_correct}, \lean{ideal\_family\_lower},
\lean{GamT\_fam} in \lean{LeanPaperA7IdealFamilyV1}; \lean{ideal\_expect\_bound},
\lean{ideal\_upper\_rate} in \lean{LeanPaperA7IdealUpperV1}; \lean{family\_policy\_lower} in
\lean{LeanPaperAFamilyLowerV1}}

The world enters the learner only as the law of its observations. Correctness is Ville's
inequality at the learner's own statistic. The upper bound comes from box-free accumulation of
the statistic on the rounds where the confidence event holds, and that event fails at round
$s$ with probability at most $s^{-4}$. The lower bound couples the stopped learner, path by
path, with a policy that queries a fresh dummy after stopping. This learner is not the
implementation measured in Section~\ref{sec:evidence}: that one has finite precision and a
bracketing oracle, and its rate is only conditional (supplement).

\begin{theorem}[What following the surrogate costs]
\label{thm:surlearner}
Let $K=4+2r$. Consider the learner with the same warm start, forcing $\frac{1}{\sqrt{s+1}}$,
tracking and stop as the ideal learner, whose oracle returns an exact optimum of the frozen
surrogate. There are two cases:
\begin{itemize}
\item the surrogate is solved at the true world and held fixed;
\item the unrestricted surrogate is re-solved at the learner's current estimate every round.
The oracle must be exact at every estimate whose extrema are the four anchors and which
retains $A$. Elsewhere it is unconstrained.
\end{itemize}
In both cases there is $\delta_{0}>0$ with
\[
\mathbb{E}\tau_{\delta}\ \ge\ \Bigl(\frac{KM}{8}\Bigr)^{2}\log^{2}(1/\delta)
\qquad(\delta<\delta_{0}),
\]
so the ratio to the ideal learner is unbounded.
\end{theorem}

\leanbox{\lean{surrogate\_optimum\_cost}, \lean{surrogate\_vs\_ideal} in
\lean{LeanPaperA6SurrogateCostV1}; \lean{est\_cost\_log\_sq}, \lean{surrogate\_learner\_vs\_ideal}
in \lean{LeanPaperA6EstCostV1}; \lean{starvation\_at} in \lean{LeanPaperA6SurrogateAtV1}}

\emph{Why.} By starvation the oracle puts nothing on the middle cells, so forcing is the only
source of middle queries. The first case is an \emph{oracle-allocation} result: the true
extrema choose the starved cells. In the second case the oracle sees only the estimate. Forcing
gives every count order $\sqrt t$, so after a deterministic time every estimate lies within
$\frac1{16}$ of the truth, except on rounds where the learner's confidence event fails. Those
rounds number at most a constant in expectation at every horizon. Along every path the middle
count is then at most a multiple of $\sqrt{\tau}$ plus that count. This is the shape the
forcing law assumes at $a=\frac12$. The Lean proofs of this theorem establish it for these two
learners directly, with the stopped change of measure at level $2\delta$ (supplement).

\emph{Which surrogate.} The estimate-time learner solves the \emph{unrestricted} frozen game at
its estimate, without centring. At the truth this game has the same value, $\frac{9}{512}$, as
the centred game of Theorem~\ref{thm:sep}. Elsewhere the two differ. Shift both arm-$A$ cells
of the truth by $u$. The unrestricted surrogate of the uniform-anchor design stays at
$\frac{9}{512}$, while the centred one is at least $\frac{9}{512}+\frac{u^{2}}{4}$. The
theorems above are about the object they name.

\leanbox{\lean{frozenGameAt\_thetaC}, \lean{frozenAt\_unif\_shift}, \lean{centred\_unif\_shift\_ge},
\lean{two\_surrogates\_differ} in \lean{LeanPaperA6TwoSurrogatesV1}}

\section{The repair law}
\label{sec:repairlaw}

Section~\ref{sec:repair} prices repair at a design and bounds what every policy must spend on
this family. This section turns the frontier $R(b)=\sup\{J_{\rm true}(w):s(w)\le b\}$ into an
operational quantity on any channel. Under a sacrifice budget $b$, the frontier is the rate
of the best stopping cost. Three quantities are tied together: the surrogate value given up,
the decision information bought with it, and the stopping time that information allows.
Multiple correct answers are not needed here: at the true world the answer is unique, so the
maximum over answers in the multiple-answer characteristic time~\cite{degenne2019} is over one
term.

\subsection{Budget classes}
Let $J_f$ be the frozen game on the channel,
$J_f(w)=\inf_{\lambda\in\mathcal L_f}\sum_qw_q\,\mathrm{KL}_q(\theta,\lambda)$, with
$\Gamma_f>0$ its value, and put $s(w)=1-J_f(w)/\Gamma_f$. The budget set
$W_b=\{w:J_f(w)\ge(1-b)\Gamma_f\}$ is convex, because $J_f$ is an infimum of linear forms.
A count vector $N$ with informative total $|N|$ \emph{meets the budget} when
$(1-b)\Gamma_f|N|\le\sum_qN_q\,\mathrm{KL}_q(\theta,\lambda)$ for every
$\lambda\in\mathcal L_f$. This is the homogeneous form of $N/|N|\in W_b$, and it needs no
division when $|N|=0$. A policy is run for $n$ rounds and queries a fresh, uninformative dummy
after it stops; its counts at the horizon are those at $\min(n,\tau)$. For realised counts
$N(\omega)$ at the horizon and expected design $\bar w=\mathbb EN/\mathbb E|N|$, consider:
\begin{itemize}
\item $\Pi_b^{\rm prefix}$: almost every history meets the budget at every prefix;
\item $\Pi_b^{\rm term}$: almost every history meets it at the horizon;
\item $\Pi_b^{\rm ev}(t_0)$: almost every history whose informative total reaches $t_0$ meets it
at the horizon;
\item $\Pi_b^{\rm exp}$: $s(\bar w)\le b$.
\end{itemize}

\begin{proposition}[The classes are nested]
\label{prop:classes}
$\Pi_b^{\rm prefix}\subseteq\Pi_b^{\rm term}\subseteq\Pi_b^{\rm exp}$ and
$\Pi_b^{\rm term}\subseteq\Pi_b^{\rm ev}(t_0)$ for every $t_0$. A policy in
$\Pi_b^{\rm ev}(t_0)$ has
\[
s(\bar w)\le b+(1-b)\,t_0/\mathbb E|N|.
\]
The constraint $\mathbb E[s(N/|N|)]\le b$ is not in this hierarchy. On the two-arm class below,
take two realised designs of lengths $1$ and $3$ with probability $\frac12$ each. Their
sacrifices average $\frac14$. Their length-weighted mean has sacrifice $\frac38$, and it carries
strictly more true information than any design within $b=\frac14$.
\end{proposition}
\begin{proof}
For each frozen alternative, integrate the pathwise inequality. Its left side integrates to
$(1-b)\Gamma_f\,\mathbb E|N|$ and its right side to $\mathbb E|N|\sum_q\bar w_q\,\mathrm{KL}_q$.
Take the infimum over $\lambda$. On the eventual class the pathwise inequality holds up to
$(1-b)\Gamma_ft_0$ on the short histories. The escape is arithmetic: the expected design is
the length-weighted mean $(5/8,3/8)$ of the realised designs $(1,0)$ and $(1/2,1/2)$, not their
plain mean.
\end{proof}
\leanbox{\lean{prefix\_terminal}, \lean{terminal\_eventual}, \lean{terminal\_expected},
\lean{eventual\_expected}, \lean{mean\_escapes} in \lean{LeanPaperAPolicyClassesV1}}

\subsection{The optimal budgeted cost}
Write $D_\delta=(1-\delta)\log(1/\delta)-\log2$.

\begin{theorem}[The repair law]
\label{thm:repairlaw}
Every channel below has finite per-query divergences $\mathrm{KL}_q(\theta,\lambda)$ at the
alternatives used, and $0<\delta\le1$.
\begin{enumerate}
\item[(i)] \emph{Every policy, any channel.} Every policy whose answer event has probability at
least $1-\delta$ at $\theta$ and at most $\delta$ at every true alternative, and which lies in
$\Pi_b^{\rm term}$, satisfies $\mathbb E|N|\cdot R(b)\ge D_\delta$. In $\Pi_b^{\rm ev}(t_0)$ it
satisfies $\mathbb E|N|\cdot R(b+(1-b)t_0/\mathbb E|N|)\ge D_\delta$. Here $R$ is the frontier
itself, a supremum over the simplex. No positivity of $\mathbb E|N|$ is assumed: with no
informative query in expectation the change of measure gives $D_\delta\le0$.
\item[(ii)] \emph{Unbounded stopping, any channel.} Let a policy stop almost surely at $\theta$,
answer wrongly at $\theta$ with probability at most $\delta$ on its stopped runs, and give the
answer at every true alternative with probability at most $\delta$ on its stopped runs. If the
expected design of its whole run meets the budget, then $\mathbb E|N(\tau)|\cdot R(b)\ge D_\delta$.
\item[(iii)] \emph{Attainment, Gaussian cells.} Let $w$ be an allocation with
$J_{\rm true}(w)>0$. Hold the oracle of the learner of Section~\ref{sec:learners} at $w$, and
keep its warm start, forcing, tracking and stop. This learner stops with a wrong answer with
probability at most $(|\mathrm{Ans}|-1)\delta$ at every world. For every
$0<\eta<J_{\rm true}(w)$ there is $A_\eta$ with
\[
\mathbb E\tau_\delta\le\frac{\log(1/\delta)}{J_{\rm true}(w)-\eta}+A_\eta
\qquad(\delta\in(0,1)).
\]
\item[(iv)] \emph{Its design enters the budget.} For every $\varepsilon>0$ there is a
deterministic $T$ such that at every $t\ge T$, on every path, the realised design dominates
$(1-\varepsilon)w$ cellwise. Hence, if $s(w)<b$, the realised design meets the budget $b$ at
every $t\ge t_0$ on every path, for some deterministic $t_0$.
\end{enumerate}
\end{theorem}
\leanbox{(i): \lean{event\_law0}, \lean{budget\_law\_alloc0}, \lean{eventual\_frontier\_law0}
in \lean{LeanPaperARepairNoPosV1}, on \lean{ebar\_alloc}, \lean{event\_law},
\lean{budget\_law\_alloc}, \lean{eventual\_frontier\_law} in \lean{LeanPaperARepairBridgeV1}, with
\lean{terminal\_eventual} and \lean{eventual\_expected} in \lean{LeanPaperAPolicyClassesV1}.
(ii): \lean{stopped\_budget\_law0}, \lean{stopped\_law\_satisfiable},
\lean{stopped\_law\_satisfiable\_two} in \lean{LeanPaperARepairNoPosV1}, on \lean{trunc\_event},
\lean{stopped\_info}, \lean{stopped\_budget\_law}, \lean{finStop\_hyps} in
\lean{LeanPaperARepairStopV1}.
(iii): \lean{oracle\_direct}, \lean{oracle\_expect\_bound}, \lean{oracle\_correct},
\lean{fix\_expect\_bound}, \lean{fix\_upper\_rate}. (iv): \lean{fix\_design\_floor},
\lean{Jval\_ge\_scale}, \lean{fix\_budget\_eventually}; all in \lean{LeanPaperABudgetUpperV1}}

In (i) the frontier bound is used only at the expected design, which is an allocation. A
bound over all nonnegative vectors would not do: $J_{\rm fr}$ and $J_{\rm true}$ are
homogeneous, so scaling a design sends its sacrifice to $-\infty$ and its true information to
$+\infty$, and no finite value of $R$ bounds it.

Part (ii) is not obtained by truncating. A policy may break the budget early and repair it
before it stops, so meeting the budget at the stop says nothing about a truncation: two queries
of arm $2$ followed by two of arm $1$ meet the two-arm budget $b=\frac12$ at the stop but not
after two rounds. The change of measure needs no budget. At a horizon where at most $\varepsilon$
of the runs at $\theta$ are unstopped, the event ``stopped with the answer'' has probability at
least $1-\delta-\varepsilon$ at $\theta$ and at most $\delta$ at every alternative. So
$(1-\delta-\varepsilon)\log(1/\delta)-\log2$ is at most the truncation's expected information,
which is at most that of the whole run, and $\varepsilon\to0$ gives $D_\delta$ itself. The budget
is then used once, at the whole run's expected design. Only the informative counts need to be
finite: after the stop the dummy is queried forever. A correct finite-horizon policy, stopped
at its horizon, meets every correctness and finiteness hypothesis of (ii). The budget is met as
well by the learners of Theorem~\ref{thm:optrate}: on Gaussian cells, the stopped fixed-design
learner at a Markov horizon, embedded as a stopping policy, meets every hypothesis of (ii) with
the budget $b+(1-b)t_0/\mathbb E|N(\tau)|$, one $t_0$ serving every $\delta<\frac18$; on two
Gaussian arms this holds for every $b\in(0,\frac12]$.

Part (iii) holds for any oracle that is eventually informative on the confidence event; the
fixed design is the case where that hypothesis is immediate. The learner reads only the
history. The design $w$ is chosen at $\theta$, as in the definition of an optimal cost
$\inf_\pi\mathbb E_\theta\tau$, so it is an oracle design. Correctness at the other worlds does
not depend on it.

\begin{theorem}[The frontier is the optimal rate]
\label{thm:optrate}
Fix a specification with Gaussian cells, its channel with a fresh dummy, a true world $\theta$
with a unique answer, frozen alternatives $\mathcal L_f$ and $\Gamma_f>0$ attained by a design.
Let $0<b\le1$ with $R(b)>0$. On the kernel policies of this channel:
\begin{itemize}
\item every family of $\delta$-correct policies in $\Pi_b^{\rm ev}(t_0)$, with one burn-in $t_0$
for all $\delta$, has $\liminf_{\delta\to0}\mathbb E|N|/\log(1/\delta)\ge1/R(b)$;
\item for every $0<\varepsilon<R(b)$ the stopped fixed-design tracking learners, run to a Markov
horizon, form such a family, with one burn-in and $\delta$-correct at every $\delta\in(0,1)$, and
$\limsup_{\delta\to0}\mathbb E|N|/\log(1/\delta)\le1/(R(b)-\varepsilon)$.
\end{itemize}
So $1/R(b)$ is the optimal budgeted rate over this class, and the learner that holds a design
inside the budget attains it. On two Gaussian arms every hypothesis holds for every
$b\in(0,\frac12]$.
\end{theorem}
\begin{proof}
\emph{The channel.} The per-query divergence of the Gaussian channel is
$(\theta_e-\lambda_e)^2/(2\sigma_e^2)$, and the dummy carries none. So the kernel information of a
design is the stream energy of its cell part, and the kernel game is the stream game.
\emph{Lower side.} This is (i) on the eventual class, with the frontier concave and hence
continuous at $b$, and $\mathbb E|N|\ge D_\delta/\Gamma_{\rm true}\to\infty$, so that the excess
$(1-b)t_0/\mathbb E|N|$ vanishes. No positivity of $\mathbb E|N|$ is assumed: a correct policy
with no informative query would give the same transcript law at $\theta$ and at an alternative.
\emph{The design.} A design within $b'<b$ and within $\varepsilon/2$ of $R(b)$ exists by
concavity. If it puts mass $p<1$ on the dummy, its cell part divided by $1-p$ has both games
multiplied by $1/(1-p)$, so its sacrifice does not increase and its true information does not
decrease.
\emph{The learner as a kernel policy.} The learner holding that design, stopped by its test and
then querying the dummy, is a kernel policy whose transcript law is the stream law of the stopped
run. At horizon $n$ its cell counts are the stream counts at some $m\le\min(n,\tau)$, and their
total is $m$. By (iv), after a deterministic burn-in the realised design meets the budget on
every path, so the policy lies in $\Pi_b^{\rm ev}(t_0)$ at every horizon, with one $t_0$ for all
$\delta$. The count vector takes countably many values, so this pathwise statement is an almost
sure one under the transcript law. At a Markov horizon the answer event of the test, run at level
$\delta/|\mathrm{Ans}|$, has probability at least $1-\delta$ at $\theta$ and at most $\delta$ at
every alternative, and $\mathbb E|N|\le\mathbb E\tau$. Part (iii) gives the rate.
\end{proof}
\leanbox{Lower side on any channel: \lean{frontier\_rate\_lower0} in
\lean{LeanPaperARepairNoPosV1}; \lean{Jt\_concave}, \lean{Jt\_le\_sum},
\lean{frontier\_rate\_lower} in \lean{LeanPaperARepairBridgeV1}, on \lean{R\_concave},
\lean{R\_continuousAt}, \lean{frontier\_lower\_rate}, \lean{near\_design} in
\lean{LeanPaperARepairRateV1}. The Gaussian channel and the class: \lean{gC}, \lean{klq\_gC\_inl},
\lean{infoAt\_gC}, \lean{Jt\_gC}, \lean{normalize\_design}, \lean{ae\_map\_counts},
\lean{real\_counts}, \lean{fix\_eventual}, \lean{fix\_etot\_le}, \lean{etot\_pos\_of\_correct},
\lean{fix\_correct}, \lean{optimal\_lower}, \lean{optimal\_upper}, \lean{optimal\_satisfiable} in
\lean{LeanPaperARepairGaussV1}. Coupling: \lean{stoppedG}, \lean{runSG\_eq},
\lean{cnt\_eq\_at\_min} in \lean{LeanPaperA7StoppedGenV1}; \lean{stream\_eq\_hist\_fresh} in
\lean{LeanPaperABridgeDummyV1}}

\subsection{The shape of repair, and a second class}
\begin{proposition}[No threshold, linear repair]
\label{prop:shape}
Let $J_f,J_t$ be concave on the simplex with $\Gamma_f=\max J_f>0$ attained.
\begin{itemize}
\item If some design has $J_t>0$, then every budget $b>0$ buys a design with $J_t>0$. There is
no threshold below which repair buys nothing.
\item With a certificate pair as in Proposition~\ref{prop:repaircertificate} and a design $v$
with $J_t(v)>0<s(v)$, the frontier is pinched on $[0,s(v)]$:
$\frac{J_t(v)}{s(v)}\,b\le R(b)\le b/\Lambda$.
\end{itemize}
\end{proposition}
\leanbox{\lean{no\_threshold}, \lean{linear\_repair} in \lean{LeanPaperARepairLawV1}}

The second class is a linear decision under Gaussian measurement, with an alternative
structure unlike the compression family. Two unit-variance arms have means $\mu_1>\mu_2$,
$\Delta=\mu_1-\mu_2$, and the decision is which mean is larger. The frozen surrogate lets only
the leader move, so $J_{\rm fr}(w)=w_1\Delta^2/2$. The true game lets both move, so
$J_{\rm true}(w)=w_1w_2\Delta^2/2$.

\begin{proposition}[Two arms]
\label{prop:twoarm}
\begin{itemize}
\item The surrogate's optimum $w=(1,0)$ has $J_{\rm true}=0$.
\item The sacrifice is $s(w)=w_2$, and $R(b)=b(1-b)\Delta^2/2$ exactly on $[0,\frac12]$.
\item The certificate pair of Proposition~\ref{prop:repaircertificate} meets the slope
$\Delta^2/2$ at $b=0$.
\item The design $(1-b,b)$ attains $R(b)$. The learner holding it stops at rate
$\limsup\mathbb E\tau_\delta/\log(1/\delta)\le2/(b(1-b)\Delta^2)$.
\item For $b'<b$, the learner holding $(1-b',b')$ meets the budget $b$ on every path after a
deterministic time.
\end{itemize}
\end{proposition}
\leanbox{\lean{jtrue2\_eq}, \lean{jproxy2\_eq}, \lean{two\_arm\_starvation},
\lean{two\_arm\_frontier}, \lean{two\_arm\_certificate} in \lean{LeanPaperARepairLawV1};
\lean{jw\_two}, \lean{jf\_two}, \lean{two\_arm\_budget\_rate}, \lean{two\_arm\_eventually} in
\lean{LeanPaperATwoArmRateV1}}

Data-sufficiency results say which directions a decision needs~\cite{bennouna2025}. Here the
needed direction is the second arm, which the surrogate's optimum never measures. The repair
law adds a price for it. Giving up a fraction $b$ of surrogate value buys decision information
$b(1-b)\Delta^2/2$. That information fixes the best stopping time,
$2\log(1/\delta)/(b(1-b)\Delta^2)$ to first order, over every policy that honours the budget
eventually. At $b=0$ the only design within the budget is the surrogate's optimum $(1,0)$, whose
true information is zero, so no policy that holds that budget can certify the decision. A
learner with forced exploration still decides, but only by spending samples outside the budget.

\section{Evidence}
\label{sec:evidence}

The evidence has three different purposes: exact certificates check finite optimisation
instances, mechanism panels reveal where a rule spends observations, and learner panels measure
finite-budget behaviour. Runtime assertions are useful consistency checks, not independent
proofs of statistical coverage or global optimality. Oracle arms read the true world and are
labelled as oracle-allocation diagnostics. For error accounting the relevant quantity is
$\Pr(\text{stop and wrong})$ over all launched runs; censoring is reported separately.
Zero observed errors on these panels is not a certification of $\delta$-correctness.

\begin{table}[t]
\caption{Evidence packages. Each script self-checks against a Lean statement.}
\label{tab:evidence}
\centering
\footnotesize
\resizebox{\columnwidth}{!}{%
\begin{tabular}{@{}llr@{}}
\toprule
Package & Script & Checks \\
\midrule
Certified separation, $r=1,2,4,8$ & \texttt{EXP\_paperA\_certified\_separation} & 6/6 \\
\quad its figure & \texttt{FIG\_paperA\_certified} & 3/3 \\
Allocation mechanism, one certificate & \texttt{EXP\_paperA\_mechanism} & 5/5 \\
\quad its figures & \texttt{FIG\_paperA\_mechanism}, \texttt{FIG\_paperA\_budget} & 7/7, 5/5 \\
Deployable rules, two-answer certificate & \texttt{EXP\_paperA\_deployable} & 6/6 \\
\quad its figure & \texttt{FIG\_paperA\_deployable} & 4/4 \\
\quad locked paired panel, $r\le4$, $M\in\{1,4\}$ & \texttt{EXP\_paperA\_deployable\_paired} & 5/5 \\
\quad its figure and paired table & \texttt{FIG\_paperA\_deployable\_paired} & 7/7 \\
Shipped learner, four-cell panel & \texttt{EXP\_paperA7\_panel} & 6/6 \\
\quad its figure & \texttt{FIG\_paperA7\_panel} & 4/4 \\
Error over launched runs & \texttt{ANALYSE\_paperA\_alpha\_launched} & 4/4 \\
Does the rule learn the answer? & \texttt{EXP\_paperA\_swap} & 7/7 \\
\bottomrule
\end{tabular}}
\end{table}

\begin{figure}[t]
\centering
\includegraphics[width=\columnwidth]{paperA_certified.pdf}
\caption{The separation, certified. A numerical true-game computation with centring and every
alternative region, certified in exact rational arithmetic (\CERT). The upper end is the dual of
the outer LP over exactly verified centred alternatives. The lower end is the minimum over all
regions of a Lagrangian dual value at a design with every weight positive.
$\Gamma_{\rm true}/\Gamma_{\rm fr}$ lies in $[0.8372,0.8374]$, $[0.7214,0.7216]$,
$[0.5664,0.5667]$ and $[0.3982,0.3982]$ at $r=1,2,4,8$ ($M=1$), inside the band of
Theorem~\ref{thm:sep}. $\Gamma_{\rm fr}=\frac{9}{512}$ is the theorem.}
\label{fig:acert}
\end{figure}

\begin{figure*}[t]
\centering
\includegraphics[width=0.92\textwidth]{paperA_mechanism.pdf}
\caption{The mechanism, from one stored panel ($M=1$, $\alpha=0.05$, 24 runs per rule and $r$,
budget $B=6\times10^{4}$); nothing was re-solved or re-sampled. Left: attained lower
certificates of the true value, up to the theorem's upper bound. Centre: the share of queries
on the middle cells. The surrogate optimum puts nothing there, and its runs reach the middles
only through forcing. Right: $\mathbb{E}[\min(\tau,B)]$. Every surrogate run is censored at $B$,
which does not show that it never stops. SCOPE: the stopping test certifies answer $A$ only,
so this is an allocation diagnostic under one certificate, not a $\delta$-correct learner
(\texttt{FIG\_paperA\_mechanism}, 7/7).}
\label{fig:amech}
\end{figure*}

\begin{figure*}[t]
\centering
\includegraphics[width=0.92\textwidth]{paperA_deployable.pdf}
\caption{Deployable rules under one two-answer certificate. Each rule is a tracking learner
with forcing $\min(1,1/\sqrt t)$, and its oracle sees only the estimate. \emph{Surrogate at
estimate} re-solves the unrestricted frozen surrogate at the estimate every round.
\emph{Geometry} ascends the certificate of the answer the estimate retains. The legend's
``geometry lower-certificate heuristic'' is this numerical ascent, not a globally certified game
optimiser. Left, centre: the
certificate rate by budget on the base world (answer $A$) and the swapped world (answer $B$),
for $r=1,2,4,8$. Right: $\mathbb{E}[\min(\tau,B)]$, $B=6\times10^{4}$, $\alpha=0.05$, 24 runs per
cell. None of the 576 launched runs stopped with a wrong answer. The 71 censored runs are all
surrogate-at-estimate runs on the base world, and the other two rules certify every run. On the
swapped world the surrogate's frozen cells are the informative ones, and it certifies every
run. SCOPE: the certificate is a lower bound on the likelihood ratio, not the exact
infimum, and no certified value of the asymptotic threshold $\delta_{0}$ is provided. The
panel shows these rules' behaviour; it does not test the $\log^{2}$ order
(\texttt{EXP\_paperA\_deployable}, 6/6; \texttt{FIG\_paperA\_deployable}, 4/4). The locked
paired panel of Figure~\ref{fig:apaired} repeats it at scale.}
\label{fig:adeploy}
\end{figure*}

\begin{figure*}[t]
\centering
\includegraphics[width=0.92\textwidth]{paperA_deployable_paired.pdf}
\caption{The locked panel. The same three rules, certificate and stop as
Figure~\ref{fig:adeploy}, with the configuration fixed and hashed before launch. There are
200 paired runs per cell: run $i$ of every rule in a cell uses the same seed. The grid is
$r\in\{1,2,4\}$, $M\in\{1,4\}$, both worlds, $B=6\times10^{4}$ and $\alpha=0.05$. Left, centre:
$\mathbb{E}[\min(\tau,B)]$ on the base and swapped worlds, with $M=1$ solid and $M=4$ dashed.
Right: the fraction certified by $B$, with exact 95\% intervals. Of 7200 launched runs, 6258
stopped, and none stopped with a wrong answer (0/2400 per rule). On the base world the
surrogate-at-estimate rule certifies 167, 41, 44, 2, 4 and 0 of 200 runs across the six
$(r,M)$ cells, while the other two rules certify every run. Its paired excess over geometry is
$1.9\times10^{4}$ to $5.4\times10^{4}$ rounds, with a bootstrap interval excluding $0$ in every
cell. Its median middle share is $0.05$, against $0.47$ for geometry. On the swapped world
every rule certifies every run, and the surrogate rule is the \emph{fastest} in every cell,
since there its frozen cells are the informative ones. SCOPE: $\min(\tau,B)$ is a restricted
mean, so where the surrogate rule is censored its mean and its excess are lower bounds. The
certificate is a lower bound and the theorem supplies no numerical $\delta_{0}$ for this panel, so this measures these
rules and does not test the $\log^{2}$ order. All 31 cell summaries of an interrupted earlier
launch are reproduced exactly (\texttt{EXP\_paperA\_deployable\_paired}, 5/5;
\texttt{FIG\_paperA\_deployable\_paired}, 7/7).}
\label{fig:apaired}
\end{figure*}

\begin{figure*}[t]
\centering
\includegraphics[width=0.92\textwidth]{paperA7_panel.pdf}
\caption{The shipped learner on a four-cell panel: base world at $M=16$ and $M=64$, a smaller
gap, and the $A/B$-swapped world. 50 runs per cell and arm, $\alpha=0.05$. Left:
$\mathbb{E}[\min(\tau,B)]$, censored runs counted. Centre: the certified lower value of the true
game at the empirical design $N/\tau$, over the game's value. Offline truth evaluates it; it
never enters the learner. Right: selection. None of the 600 launched runs stopped with a wrong
answer; one learner run on the swapped world was censored. The panel exposes an algorithmic
weakness. The learner's kept set never narrowed to a single answer in any cell, so its
allocation follows the tie-break answer. On the swapped world that answer is the wrong one, the
learner selected the true answer in none of its refreshes, and it cost $3801$ queries against
$1109$ for the oracle-allocation diagnostic and $1679$ for uniform. Its stop was still correct.
SCOPE: cost estimation on a finite panel, not error certification
(\texttt{EXP\_paperA7\_panel}, 6/6; \texttt{FIG\_paperA7\_panel}, 4/4).}
\label{fig:apanel}
\end{figure*}

\paragraph{What the panels establish}
The locked panel isolates a directional failure, not universal superiority of a new allocator.
On the base world at $(r,M)=(4,4)$, the surrogate rule reaches the $60000$ cap in all 200 runs,
whereas geometry and uniform have restricted means $13682.25$ and $15622$.
At $(1,1)$, geometry and uniform are nearly tied ($1737.5$ and $1732.25$).
On the swapped world the surrogate rule is fastest in all six parameter cells. These reversals
are informative: the cost is tied to which directions the frozen geometry excludes.
The implemented geometry rule uses a numerical ascent of a lower certificate; it is not a
globally certified maximiser of the true game. Likewise, the shipped learner's failure to
resolve its answer-selection tie on the swapped world is an implementation limitation that
the ideal learner theorem does not excuse.

\section{Limitations}
\label{sec:limits}

\begin{enumerate}
\item[1)] \textbf{One construction for the rates.} Theorems~\ref{thm:sep}--\ref{thm:surlearner}
are about the centred construction. Both directions of mispricing are proved after the sup over
designs on a second class, but its learners and rates are not studied here.
\item[2)] \textbf{Spreading the budget, not reducing it.} Theorem~\ref{thm:dilute} spreads a
paid budget over any run length by padding after the decision is made. The padded queries carry
no information for the decision, so the spreading is an accounting fact. It does not show that
interleaving surrogate-optimal queries \emph{during} the run lowers the budget below $N_{0}$.
At a design the constant is exact: the frontier's slope at zero is $\frac{1}{8rM}$
(Theorem~\ref{thm:frontier}). At the level of policies the lower bound $8rMD$ is not matched.
The ideal learner is time-optimal, not budget-optimal, and its budget exceeds $8rMD$ by at most
the factor $\frac{32}{9}+\frac{64}{9rM}$ of Theorem~\ref{thm:sacrifice}. A learner whose budget is
$8rM\log(1/\delta)(1+o(1))$ is not constructed here; Theorem~\ref{thm:optrate} is not
instantiated on this family.
\item[3)] \textbf{The implemented learner.} The learner of Figure~\ref{fig:apanel} has finite
precision and a bracketing oracle. Its rate is conditional on named accumulation and
confidence hypotheses (supplement), and its selection step did not narrow on the panel.
\item[4)] \textbf{Stopping.} The sacrifice law and the forcing law hold for unbounded stopping at
level $2\delta$ (Theorems~\ref{thm:sacstop}, \ref{thm:forcing}); the repair law at the stop holds
at level $\delta$ (Theorem~\ref{thm:repairlaw}(ii)). The forcing law for stopping asks $\varphi$
to be nondecreasing.
\item[5)] \textbf{The repair law.} The attaining design of Theorem~\ref{thm:optrate} is chosen
at the true world, as in the definition of an optimal cost. A deployable learner that re-solves
the budgeted game at its estimate needs continuity of the constrained maximiser and is not
treated. The optimal rate is proved on Gaussian cells; on other channels only the lower side is.
With several correct answers only the bound with $(1-\delta)/|\text{correct}|$ is proved, for a
policy with a positive expected informative count (\lean{budget\_law\_multi}); the exact
multi-answer limit is not claimed.
\item[6)] \textbf{Certificates at finitely many points.} The separation of
Figure~\ref{fig:acert} is certified exactly at $r=1,2,4,8$. The two fixed-design families of the
supplement are certified exactly, in all 81 regions, at six values of $M$. Their closed forms for
general $M$ are read off those points and not proved.
\end{enumerate}

\section{Conclusion}

The relevant failure is not merely that two design objectives assign different numbers.
On a realizable compression family, every exact and near-optimal frozen design loses the
information needed by the actual decision. A dual-margin certificate turns that geometric
failure into a design tradeoff and then into a policy-level budget. The family makes its
scale explicit: for $rM\ge1$, repair costs $\Theta(rM\log(1/\delta))$ query-units of surrogate
value, and at a design the exchange rate $1/(8rM)$ is exact. With square-root corrective exploration, the specified surrogate learner pays
$\Omega(\log^2(1/\delta))$, while the ideal true-geometry learner attains its characteristic
time. On any channel the frontier bounds the rate under a sacrifice budget, and on Gaussian
cells a tracking learner attains it within the same class, so the price of repair is paid in
stopping time as well as in surrogate value. The finite panels
support the mechanism and expose implementation weaknesses; they
do not establish the asymptotic rates or universal superiority over uniform allocation.

\section*{Companion material}
The supplement gives the proof architecture of every theorem, the supporting Lean declarations,
the conditional status of the implemented learner's rate, the heteroscedastic families, the
repair criterion and the confidence-sequence machinery, with each superseded statement marked.
The claim map records each headline's model, declaration, scope and evidence.

\begin{thebibliography}{00}

\bibitem{degenne2019} R. Degenne and W. M. Koolen, ``Pure exploration with multiple correct
answers,'' in \emph{Advances in Neural Information Processing Systems 32}, 2019,
pp.~14591--14600.

\bibitem{wangyou2026} Q. Wang and W. You, ``Ranking-and-selection with multiple correct
answers and non-answerable estimates,'' arXiv:2606.21889, 2026.

\bibitem{bennouna2026} O. Bennouna, A. Bennouna, S. Amin, and A. Ozdaglar, ``Data
informativeness in linear optimization under uncertainty,'' arXiv:2602.15365, 2026.

\bibitem{bennouna2025} O. Bennouna, A. Bennouna, S. Amin, and A. Ozdaglar, ``What data enables
optimal decisions? An exact characterization for linear optimization,'' arXiv:2505.21692, 2025.

\bibitem{garivier2016} A. Garivier and E. Kaufmann, ``Optimal best arm identification with
fixed confidence,'' in \emph{Proc. 29th Conf. Learning Theory (COLT)}, PMLR 49, 2016,
pp.~998--1027.

\bibitem{kaufmann2016} E. Kaufmann, O. Capp\'e, and A. Garivier, ``On the complexity of
best-arm identification in multi-armed bandit models,'' \emph{J. Mach. Learn. Res.}, vol.~17,
no.~1, pp.~1--42, 2016.

\bibitem{jourdan2021} M. Jourdan, M. Mutn\'y, J. Kirschner, and A. Krause, ``Efficient pure
exploration for combinatorial bandits with semi-bandit feedback,'' in \emph{Proc. 32nd Int.
Conf. Algorithmic Learning Theory (ALT)}, PMLR 132, 2021, pp.~805--849.

\bibitem{shukla2024} A. Shukla and D. Basu, ``Preference-based pure exploration,''
arXiv:2412.02988, 2024.

\bibitem{yang2024} J. Yang, V. Y. F. Tan, and T. Jin, ``Best arm identification with minimal
regret,'' arXiv:2409.18909, 2024.

\bibitem{huang2026} Z. Huang, ``Self-certification of representation adequacy: Sequential
certification at minimum task loss,'' arXiv:2608.02267, 2026.

\bibitem{schutte2026} N. Schutte, G. Veviurko, K. Postek, and N. Yorke-Smith, ``Sufficient
decision proxies for decision-focused learning,'' arXiv:2505.03953, 2026.

\bibitem{overstall2022} A. M. Overstall, ``Properties of using Fisher information gain for
Bayesian design of experiments,'' arXiv:2003.07315, 2021.

\bibitem{filstroff2021} L. Filstroff, I. Sundin, P. Mikkola, A. Tiulpin, J. Kylm\"{a}oja,
and S. Kaski, ``Targeted active learning for Bayesian decision-making,'' arXiv:2106.04193,
2021.

\bibitem{tang2026} R. Tang, S. J. Sloman, and S. Kaski, ``Representative, informative, and
de-amplifying: Requirements for robust Bayesian active learning under model misspecification,''
in \emph{Proc. Int. Conf. Artificial Intelligence and Statistics (AISTATS)}, 2026,
arXiv:2506.07805.

\bibitem{wan2026} B. Wan, M. Liu, P. Grigas, and Z.-J. M. Shen, ``Decision-focused
sequential experimental design: A directional uncertainty-guided approach,'' arXiv:2602.05340,
2026.

\bibitem{mathlib2020} The mathlib Community, ``The Lean mathematical library,'' in
\emph{Proc. 9th ACM SIGPLAN Int. Conf. Certified Programs and Proofs (CPP)}, 2020,
pp.~367--381.

\bibitem{demoura2021} L. de Moura and S. Ullrich, ``The Lean 4 theorem prover and
programming language,'' in \emph{Proc. 28th Int. Conf. Automated Deduction (CADE)}, 2021,
pp.~625--635.

\end{thebibliography}

\end{document}
